\chapter*{Abstract}
Monocular 3D object detection, a fundamental capability for autonomous driving and robotics, aims to localize objects in 3D space from a single RGB image. The task remains challenging due to the inherent ambiguity of recovering 3D geometry from 2D observations. This thesis presents two complementary contributions that advance monocular 3D detection along distinct axes: the efficiency of feature extraction and the geometric consistency of the detection objective.

\textit{In the first work}, we observe that Transformer-based monocular 3D detectors incur substantial overhead due to the quadratic cost of self-attention. We propose MonoDSSMs, the first monocular 3D detector to leverage the Mamba architecture for efficient feature extraction. It introduces a Bidirectional Mamba2 (BiMamba2) module that captures global context with linear complexity, a CrossMamba module for fusing depth-aware and context-aware features, and a multi-scale depth enhancement strategy. On KITTI, MonoDSSMs matches its baseline's detection quality while using 2.2x fewer parameters and running 1.28x faster, and generalizes across modern Vision Mamba architectures with both fixed-scan and learnable-scan strategies as plug-and-play replacements.

\textit{In the second work}, we address a critical limitation of DETR-like monocular 3D detectors: the exclusion of 3D attributes from bipartite matching, which causes high-quality 3D predictions to be suppressed by 2D-only criteria. We propose Mono3DV, a Transformer-based framework with three innovations: a 3D-Aware Bipartite Matching strategy that incorporates 3D geometry into the matching cost, a 3D-DeNoising scheme that stabilizes training under 3D criteria, and a Variational Query DeNoising mechanism that overcomes the gradient vanishing issue of conventional denoising. Without external data, Mono3DV achieves state-of-the-art results on KITTI and Waymo, and delivers consistent gains as a plug-and-play module for multi-view detectors on nuScenes.